\subsection{A coefficient-scaled bar image gives a natural acyclic disk envelope}
\label{arch136:scaled-bar}

\paragraph{ARCH-136-SCALED-BAR.}
Let \(V\) be a free abelian group with a basis of even form degrees at
least two and positive weights. Write \(b_v\) for its generators and
\(a_v\) for a copy one form degree lower, of the same weight, and put
\[
 Q=\Gamma_{\mathbb Z}(V),\qquad
 D=\Gamma_{\mathbb Z}(V)\otimes
       \Lambda_{\mathbb Z}(s^{-1}V),\qquad
 d a_v=b_v,\quad d b_v=0.
\]
Give both \(a_v\) and \(b_v\) auxiliary arity one.
There is a strict differential PD \(S=\Gamma_{\mathbb Z}(T)\)-subalgebra
\[
 E(V)\ \subseteq\ S\otimes_{\mathbb Z}D
\]
which contains \(T a_v,T b_v\), is natural for graded,
weight-preserving linear maps of \(V\), and has a weight-preserving
\(S\)-linear contraction onto \(S\). Thus
\[
 H^0(E(V))=S,\qquad H^q(E(V))=0\quad(q>0)
\]
in every weight. This is a disk-envelope statement. It does not assert
that an arbitrary simplicial partial source embeds in such an envelope,
or that arbitrary coefficient-tail cycles are generated by its scaled
outputs.

\paragraph{Integral Euler operator.}
Define the odd derivation \(h\) on \(D\) by
\[
 h(b_v)=a_v,\qquad h(a_v)=0,\qquad
 h(\gamma_r(b_v))=a_v\gamma_{r-1}(b_v).
\]
This is integral. On a monomial it is the sum of divided-power partial
derivatives multiplied by the corresponding exterior generators.
The diagonal terms in \(h^2\) vanish because \(a_v^2=0\); distinct
terms cancel in pairs because the divided-power partial derivatives
commute and the exterior generators anticommute. Hence \(h^2=0\).
The commutator \(dh+hd\) is the arity Euler operator:
\[
 (dh+hd)|_{D_r}=r\,\mathrm{id}.
 \tag{SB.1}
\]
In particular \(H=h/r\) is a contraction of the positive-arity part
of \(D_{\mathbb Q}\).

\paragraph{Explicit bar comparison and denominator bound.}
Use the normalized one-sided bar resolution
\(\mathsf B(Q)=Q\otimes T(s^{-1}\bar Q)\), with its shuffle product.
For a word \(u_0[u_1|\cdots|u_k]\), each \(u_j\) is homogeneous in
arity \(r_j\), and \(r_j\ge1\) for \(j\ge1\).
All \(u_j\) have even form degree. For \(k\ge1\), put
\[
 R_j=r_j+\cdots+r_k,\qquad P=R_1R_2\cdots R_k,
\]
and define over \(\mathbb Q\)
\[
 \rho\bigl(u_0[u_1|\cdots|u_k]\bigr)
   =\frac{u_0h(u_1)\cdots h(u_k)}{P}.
 \tag{SB.2}
\]
For \(k=0\), set \(\rho(u_0)=u_0\). Formula (SB.2) is extended
linearly. The integers \(R_1,\ldots,R_k\) are distinct elements of
\(\{1,\ldots,R_1\}\), so \(P\mid R_1!\).

Let \(r=r_0+\cdots+r_k\). Define
\[
 \Phi\bigl(u_0[u_1|\cdots|u_k]\bigr)=
       T^r\,\rho\bigl(u_0[u_1|\cdots|u_k]\bigr).
 \tag{SB.3}
\]
Because \(P\mid r!\), the coefficient \(T^r/P=(r!/P)T^{[r]}\)
is integral. Thus (SB.3) defines an integral \(S\)-linear map
\(S\otimes\mathsf B(Q)\to S\otimes D\).
To make it weight-preserving, give the source generator \(b_v\) weight
\(\operatorname{wt}(b_v)+1\); bar suspension lowers form degree by
one and preserves this new weight. The target has its original
weights and \(\operatorname{wt}(T)=1\). A word of positive arity \(r\)
has form degree at least \(r\), so no negative form degrees arise.

\paragraph{Chain-map verification.}
For the chosen bar convention,
\[
\begin{split}
 d_{\mathrm{bar}}\bigl(u_0[u_1|\cdots|u_k]\bigr)
  ={}&u_0u_1[u_2|\cdots|u_k]\\
   &+\sum_{i=1}^{k-1}(-1)^i
       u_0[u_1|\cdots|u_i u_{i+1}|\cdots|u_k].
\end{split}
 \tag{SB.4}
\]
There is no final augmentation term because \(u_k\in\bar Q\).
Since \(d h(u_j)=r_j u_j\), differentiating (SB.2) gives terms
with coefficient \((-1)^{j-1}r_j/P\).
For a merged pair \(u_i u_{i+1}\), its denominator is \(P/R_{i+1}\).
Using
\(h(u_i u_{i+1})=h(u_i)u_{i+1}+u_i h(u_{i+1})\),
the coefficient of the term containing \(u_j\) and every other
\(h(u_\ell)\) is
\[
 \frac{R_1-R_2}{P}\quad(j=1),\qquad
 \frac{(-1)^{j-1}(R_j-R_{j+1})}{P}\quad(1<j<k),
 \qquad \frac{(-1)^{k-1}R_k}{P}\quad(j=k).
\]
These are the required coefficients. The case \(k=1\) is
\(d(h(u_1)/r_1)=u_1\), and \(k=0\) is immediate.
Thus \(\rho\), and hence \(\Phi\), is a chain map.

\paragraph{Shuffle multiplication and PD operations.}
For positive integers \(r_1,\ldots,r_k\), the reciprocal denominator
in (SB.2) is the iterated integral
\[
 \frac1{R_1\cdots R_k}
 =
 \int_{0\le t_k\le\cdots\le t_1\le1}
       t_1^{r_1-1}\cdots t_k^{r_k-1}\,dt_1\cdots dt_k.
 \tag{SB.5}
\]
The product of two such ordered simplices is partitioned, apart from
their common boundaries, by the shuffles of the two orderings.
Thus the product of the two reciprocal denominators is the sum of
the reciprocal denominators for all shuffles. In the bar shuffle
product every letter is odd after suspension. The shuffle sign
cancels the identical sign from reordering the odd factors
\(h(u_j)\) in the target. Equation (SB.5) therefore proves that
\(\rho\) is multiplicative. Scaling an arity-\(r\) element by \(T^r\)
also respects multiplication, so \(\Phi\) is multiplicative.

The shuffle bar algebra carries its classical strict PD structure;
see Cartan, \emph{Puissances divis\'ees}, Expos\'e~7,
Theorem~1, pp.~7-01--7-02. Over \(\mathbb Q\), divided powers on
even elements are \(x^n/n!\), so the multiplicative map \(\rho\),
and its coefficient-scaled map, preserve them.
Both the bar algebra and \(S\otimes D\) are torsion-free as abelian
groups; the preceding integral image calculation therefore proves
the same identities integrally. Thus the image \(E(V)=\operatorname{im}\Phi\)
is a strict differential PD \(S\)-subalgebra.

\paragraph{Contraction of the image.}
The operator \(H=h/r\) is \(S\)-linear on target arity \(r>0\);
the coefficient algebra \(S\) has auxiliary arity zero.
For a homogeneous bar word as above,
\[
 H\,\rho\bigl(u_0[u_1|\cdots|u_k]\bigr)
 =
 \rho\bigl(1[u_0|u_1|\cdots|u_k]\bigr)
 \quad\text{if }r_0>0.
 \tag{SB.6}
\]
Indeed \(h^2=0\), so only the derivative of \(u_0\) contributes,
and the new denominator is \(rP\). If \(r_0=0\), the left side
is zero. The insertion in (SB.6) preserves total arity, so the
same identity holds after the \(T^r\)-scaling in (SB.3).
Consequently \(H\) preserves \(E(V)\) integrally, independently of
a chosen presentation of an image element. By (SB.1),
\(dH+Hd=\mathrm{id}\) in positive arity, and it is zero on the
coefficient copy of \(S\). This proves the claimed contraction.
All operators and formulas commute with linear maps of \(V\),
which proves the stated naturality. \(\square\)

\paragraph{One-generator form.}
For one odd generator \(a\), \(da=b\), the image has the \(S\)-basis
\(1,a_r,b_r\), \(r\ge1\), where
\[
 b_r=T^r\gamma_r(b),\qquad
 a_r=\frac{T^r}{r}a\gamma_{r-1}(b)
     =(r-1)!T^{[r]}a\gamma_{r-1}(b),
 \qquad da_r=b_r.
\]
Words of length at least two map to zero because \(a^2=0\).
The products are
\[
 b_m b_n=\binom{m+n}{m}b_{m+n},\qquad
 b_m a_n=\binom{m+n}{m}a_{m+n},\qquad
 a_m a_n=0.
\]
The displayed basis elements are independent over \(S\), since
multiplication by each nonzero \(T^r/r\) is injective in \(S\).
This verifies the proposed one-disk envelope and its coherent
many-generator extension.

\paragraph{What the special conormal does and does not generate.}
In the notation of (CON.8), with \(N\ge2\) and
\(R_0=\mathbb Z[1/(2N)!]\), one has
\[
 \bigl(I_0/I_0^{[2]}\bigr)/
 A_{0,+}\bigl(I_0/I_0^{[2]}\bigr)
 \cong R_0 e_N\oplus\bigoplus_{p>2N}\mathbb F_p v_p.
 \tag{SB.7}
\]
Hence the bottom class \(v_p\) at weight \(p\) is not a positive
coefficient multiple of any lower conormal class.
The conormal calculation therefore identifies necessary
generators; it does not supply the missing assertion that arbitrary
coefficient-tail relations are generated by \(T^r\)-scaled outputs.
Nor is the stated naturality a claim about arbitrary nonlinear
simplicial attaching maps. Both bridges remain separate requirements
for applying this envelope to the integral common-source problem.

